%% =====================================================================================
\section{强子化修正}\label{sec:jets-hadronization}

推力一类事件形状的强子化修正是 $\Lambda/Q$ 量级（$\Lambda\sim\Lqcd$）而非“扭度”直觉的 $\Lambda^2/Q^2$（误解 M10）：在推力的因子化（\cref{thm:scet-thrust}，全阶论证，[A]）中，
领头幂次 $\tau=(p_{\rm L}^2+p_{\rm R}^2)/Q^2+k/Q$，软动量 $k$ 以一次方进入，$k\sim\Lambda$ 的涨落使 $\tau$ 移动 $\sim\Lambda/Q$。设在 $Q\tau\gg\Lambda$ 的尾部，
软函数是微扰软函数与归一化形状函数 $F$ 的卷积，$S_T(k)=\int\dd k'\,S_T^{\rm PT}(k-k')F(k')$，$\Omega_1\equiv\int\dd k'\,k'F(k')\sim\Lambda$。代入\cref{thm:scet-thrust}的卷积（令 $k=k''+k'$），
\begin{equation}\label{eq:jets-shift-derivation}
\frac{\dd\sigma}{\dd\tau}(\tau)=\int\dd k'\,F(k')\,\frac{\dd\sigma_{\rm PT}}{\dd\tau}\Big(\tau-\frac{k'}{Q}\Big)
=\frac{\dd\sigma_{\rm PT}}{\dd\tau}\Big(\tau-\frac{\Omega_1}{Q}\Big)\Big[1+\order{\frac{\Lambda^2}{Q^2\tau^2}}\Big],
\end{equation}
第二个等号是在 $\tau-\Omega_1/Q$ 处的 Taylor 展开：一次项按 $\Omega_1$ 的定义为零，二阶项含 $\langle(k'-\Omega_1)^2\rangle_F/Q^2$ 与 $\dd\sigma_{\rm PT}/\dd\tau\sim\ln\tau/\tau$ 的二阶导数，
相对大小为 $\order{\Lambda^2/(Q\tau)^2}$。尾部的领头效应于是整体平移 $\Omega_1/Q$，相对大小 $\Omega_1/(Q\tau)$ 随 $\tau$ 减小而增大，$Q\tau\sim\Lambda$ 时须用整个形状函数\cite{Korchemsky:1999kt,Bauer:2003di}。
其他事件形状 $e$ 同样平移 $c_e\Omega_1/Q$；问题在于卷积结构是否成立、各事件形状的 $\Omega_1$ 是否相同——后者即下面的猜想。

\begin{conjecture}[强子化修正的普适性 \tagC]\label{conj:jets-hadronization}
对一族两喷注事件形状 $e$（推力 $\tau$ 等），在 $\Lambda/Q\ll e\ll1$ 的尾部，领头的强子化修正是分布的平移 $e\to e-c_e\,\Omega_1/Q$：$c_e$ 是由观测量决定的可计算系数，
$\Omega_1\sim\Lambda$ 是对整族事件形状\emph{相同}的非微扰参数；相应地，平均值 $\langle e\rangle$ 受到 $c_e\Omega_1/Q$ 的修正。在 Dokshitzer 与 Webber 的模型中，$\Omega_1$ 由有效耦合在红外区的积分给出\cite{DokshitzerWebber1995}。
非微扰地位：\cref{asm:intro-bridge}（MA1、MA2）。本猜想以其 (Br2)——全部 $\Lambda$ 敏感性只出现在软函数中——为前提，而它所断言的普适性超出该假设。
\uses{thm:scet-thrust,asm:intro-bridge}
\end{conjecture}

\paragraph{地位与缺口。}
在\cref{thm:scet-thrust}的框架内，全部 $\Lambda$ 敏感性都在推力软函数——Wilson 线乘积的真空矩阵元——之中，$Q\tau\sim\Lambda$ 时它是非微扰的形状函数
\cite{Korchemsky:1994is,Korchemsky:1999kt,Bauer:2003di,Lee:2006nr,Bauer:2008dt}。据这些文献的摘要：形状函数与 $Q$ 无关，组织全部 $1/(eQ)^n$ 型修正，但对不同事件形状\emph{并不相同}，
只携带能量—动量张量的普适矩阵元的信息\cite{Korchemsky:1999kt}；平均值的领头幂次修正由单个动量流算符与 Wilson 线的关联决定\cite{Lee:2006nr}；事件形状分布可写成有效理论的真空矩阵元，
而不必假定强子末态按部分子的共线与软态分解\cite{Bauer:2008dt}。Dokshitzer--Webber 模型中的普适性建立在有效耦合于红外区行为正常这一假设之上\cite{DokshitzerWebber1995}。
因此“整族事件形状共享同一个 $\Omega_1$”是猜想（[C]）。重整化子提供了 $\Lambda/Q$ 修正的微扰动机：大 $\beta_0$ 近似中事件形状的红外重整化子给出 $1/Q$ 歧义
（\cref{prop:power-renormalon-examples}，部分结果，[P]；前向指引）。升级途径：对一族事件形状给出 $\Omega_1$ 的非微扰定义并证明其普适性（例如借助格点或可解模型），这仍以\cref{asm:intro-bridge}为前提。

\begin{tcolorbox}[enhanced,breakable,title=常见误解（M8、M9、M10）,fonttitle=\bfseries,
  colback=gray!3!white,colframe=gray!55!black,boxrule=0.6pt,arc=2pt,
  left=5pt,right=5pt,top=3pt,bottom=3pt,before skip=6pt,after skip=6pt]
\textbf{M8} “KLN 定理解释了强子对撞截面为何有限。”——初态共线奇异性不相消，而被吸收进部分子分布（如\cref{thm:dy-factorization}）；含初态简并求和的 KLN（\cref{thm:ir-kln}）不对应强子束流的截面；
有色初态的 Bloch--Nordsieck 相消在两圈失效（\cref{prop:ir-bn-failure}）。\par
\textbf{M9} “红外安全保证了因子化与微扰控制。”——红外安全只处理末态奇异性（\cref{cor:jets-ir-finite}），有入射强子时仍须因子化；它仍可有非全局对数（\cref{prop:jets-ngl}）、
超领头对数与 Glauber 敏感性（\cref{prop:jets-sll}）及 $\order{\Lambda/Q}$ 的强子化修正（\cref{conj:jets-hadronization}）。\par
\textbf{M10} “幂次修正总是 $\order{\Lambda^2/Q^2}$。”——事件形状的强子化修正是 $\order{\Lambda/Q}$（\cref{sec:jets-hadronization}），锥喷注截面的则按 $1/(Q\delta)$（据\cite{Korchemsky:1994is}的摘要）；见\cref{ch:power}。
\end{tcolorbox}
